% ##############################################################################
\section{Open questions}
\label{sec:open_questions}

The design in Sections~\ref{sec:noesis_market} and \ref{sec:noesis_server}
is a starting point rather than a settled protocol. This section
consolidates the open questions raised throughout the paper, grouped by
which component they primarily affect.

% ==============================================================================
\subsection{Market design questions}
\label{sec:open_questions_market}

\begin{enumerate}
  \item \textbf{Task unit}: which definition of the task unit
    (Definition~\ref{def:task}), raw tokens, wall-clock compute, or a
    benchmark-normalized task-equivalent, gives the best cross-provider
    comparability, and how sensitive is the auction outcome to this
    choice?

  \item \textbf{Auction frequency and mechanism}: is a uniform-price call
    auction on a fixed $T$-minute cadence
    (Section~\ref{sec:rounds}) the right mechanism at this scale, or would
    a continuous double auction serve latency-sensitive buyers better,
    mirroring the tension between frequent batch auctions and continuous
    limit order books \cite{budish2015}? A hybrid design, a fast spot
    market layered on top of the periodic batch auction, mirroring the
    day-ahead/real-time separation in electricity markets
    (Section~\ref{sec:related_commodity_markets}), is a natural candidate
    we have not evaluated.

  \item \textbf{Compatibility versus scoring}: Definition~\ref{def:compatibility}
    treats latency and reliability as hard eligibility constraints. Would
    a scored compatibility measure, admitting partial matches at a
    discount, recover otherwise-unmatched volume without requiring an
    explicit exchange rate between quality and price?

  \item \textbf{Pricing denomination}: should $\priceattr$ be denominated
    in a real-world currency or in a synthetic task-credit, and how much
    does that choice affect adoption and regulatory exposure?

  \item \textbf{Resistance to gaming without gatekeeping}: how far can
    Section~\ref{sec:mechanism_design_risks}'s ex post reputation penalty
    substitute for ex ante seller vetting, without the market drifting
    toward the kind of slow onboarding or reputation gate that would
    defeat the goal of an open market?

  \item \textbf{On-chain settlement and anonymity}: Section~\ref{sec:decentralized_extension}
    sketches a decentralized, blind-bid variant of \NoesisMarket; whether
    that design's gas costs, redundancy defaults, and stake asymmetries
    (Section~\ref{sec:decentralized_open_questions}) are compatible with
    the auction frequency of Section~\ref{sec:rounds} remains open.
\end{enumerate}

% ==============================================================================
\subsection{Server design questions}
\label{sec:open_questions_server}

\begin{enumerate}
  \item \textbf{Where the value actually lies}: is a simple pass-through
    logger, without any routing intelligence, already useful for building
    a dataset for downstream distillation, or does routing quality itself
    materially affect the diversity and usefulness of the collected data?

  \item \textbf{Net savings from routing}: how cheaply can the difficulty
    estimator of Equation~\eqref{eq:difficulty_estimator} be trained and
    served while still satisfying the saving condition of
    Equation~\eqref{eq:routing_saving_condition}, and what fraction of
    real traffic can be served by a small model with no measurable
    quality loss?

  \item \textbf{Routing versus fusion under a fixed budget}: for a fixed
    per-request budget, is it preferable to route to one well-chosen
    model (Section~\ref{sec:routing}) or to query several cheaper models
    and combine their answers (Section~\ref{sec:fusion})?

  \item \textbf{Privacy and data handling}: what safeguards, PII
    scrubbing, opt-in versus opt-out consent, retention limits, are
    required before \NoesisServer{} logs real prompt/response pairs at
    scale (Section~\ref{sec:gateway_architecture})?

  \item \textbf{Distillation quality}: can a model distilled from the
    logged corpus (Section~\ref{sec:distillation}) match the quality of
    the routed ``best tier per request'' policy at a fraction of the
    inference cost, or does distillation only recover the easy-request
    tail?

  \item \textbf{Attribution reliability}: the hypothesis test of
    Equation~\eqref{eq:reliability_lower_bound} distinguishes sampling
    noise from a true reliability shortfall, but it still relies on
    $\hat c(x)$ and $\hat\ell(x)$ being accurate per-request estimates of
    delivered capability and latency; how reliable can those per-request
    estimates be made, and how does verifier error propagate into false
    or missed violations reported to \NoesisMarket?
\end{enumerate}

% ==============================================================================
\subsection{Cross-cutting questions}
\label{sec:open_questions_cross}

The reputation and pricing feedback loop of Section~\ref{sec:reputation}
is only as trustworthy as the fulfillment monitoring that feeds it
(Section~\ref{sec:fulfillment_monitoring}). Resolving the server-side
attribution-reliability question above is therefore a prerequisite, not
an independent improvement, for the market-side reputation mechanism to
have the intended deterrent effect on capability misrepresentation and
under-delivery.

\textbf{Interface contracts for pluggable components}: Section~\ref{sec:modularity}
treats the matching engine, capability measurement, reputation and
feedback, answer fusion, and pricing dissemination as independently
replaceable components, but this paper only sketches each interface in
prose. Making a component genuinely swappable, e.g., trying a continuous
double auction in place of the call auction of Section~\ref{sec:noesis_market}
without touching \NoesisServer{}, requires pinning down the exact inputs,
outputs, and invariants each interface guarantees, which is left as
future specification work rather than resolved here.

% ==============================================================================
\subsection{Open questions specific to the decentralized design}
\label{sec:decentralized_open_questions}

\begin{enumerate}
  \item \textbf{Bid and price discovery for a new market}: with no trading
    history, neither a buyer nor a seller has a reference price to anchor
    a bid or an ask to; whether the first rounds of a new tier bucket need
    a different mechanism, e.g., a wider price band or a subsidized
    bootstrap period, than the steady-state design above is unresolved.

  \item \textbf{Reveal-window griefing}: a blind bid
    (Definition~\ref{def:blind_bid}) is visible to on-chain observers once
    revealed; a party monitoring the mempool during the reveal window can
    still submit a competing bid or ask before the window closes, unless
    the window is bounded to a single block, which in turn bounds how many
    participants can feasibly reveal in time.

  \item \textbf{Commitment without capability}: nothing in
    Definition~\ref{def:blind_bid} stops a seller from committing an ask
    it cannot actually fulfill within $\lmaxattr$, since the stake
    $\sigma_\alpha$ is only forfeited after the fact; whether a smaller,
    non-refundable commitment fee at bid time is needed to discourage
    speculative asks is open.

  \item \textbf{Stake sizing and asymmetry}: whether the stake
    $\sigma_\alpha$ required of a seller (Section~\ref{sec:stake_slashing})
    should be set unilaterally by the buyer's bid, which risks the buyer
    demanding excessive collateral, or by a market-wide rule tied to the
    contract's price, and whether the buyer should also be required to
    post a bond so that an under-collateralized buyer cannot costlessly
    request contracts it does not intend to honor payment for.

  \item \textbf{Sybil resistance}: whether stake-and-slash
    (Section~\ref{sec:stake_slashing}) is by itself a sufficient deterrent
    against a seller re-entering under a new address after being slashed,
    or whether a Sybil-resistant identity layer is still needed on top, in
    tension with the anonymity goal of Section~\ref{sec:decentralized_motivation}.

  \item \textbf{Redundancy defaults}: whether $k_{\text{red}}$
    (Section~\ref{sec:decentralized_redundancy}) should default to 1 and
    be raised only at the buyer's request and expense, or whether the
    market itself should raise it automatically in thin, collusion-prone
    buckets.

  \item \textbf{Verifiable optimality}: whether a succinct proof of
    clearing-problem optimality, not merely feasibility
    (Section~\ref{sec:decentralized_matching_engine}), is practical to
    generate and verify at the scale of a real tier bucket, or whether the
    design must settle for verified feasibility and accept a centralized
    operator's word on optimality.
\end{enumerate}
